\section{Parameter-preserving fragments with arbitrary heads}
\label{m:sec:certificates}
Polynomial graphs were needed for completely unrestricted arrangements of
polynomial heads. Different arrangements permit arbitrary fixed heads
without constructing their graphs. This section isolates the spacing fact
used by all of them.

\begin{lemma}[Sparse endpoint comparison]\label{m:lem:endpoints}
For a fixed nonzero limit $\lambda$, $J_\lambda$ is strictly increasing.
Its distinct norm inputs have distinct terminal endpoints, ordered in the
same way. For any two distinct endpoints $Z<W$ one has
\begin{equation}\label{m:eq:sparse}
 W\ge2^{Z+1}.
\end{equation}
Every value with endpoint $Z$ lies strictly between $Z$ and $2Z$ and below
the next value of $f$.
\end{lemma}
\begin{proof}
Strict increase is controlled strictness from Lemma~\ref{f:lem:collapse-estimates}: when $t<u$, the smaller ordinal has norm
$t\le f(u)$, so the controlled comparison premise holds. Write
$J_\lambda(t)=f^L(t)+L$ and $J_\lambda(u)=f^M(u)+M$.
If their endpoints were equal, injectivity of $f$ would force $L>M$;
otherwise cancelling $L$ iterates would give $t\ge u$. The equal endpoints
and $L>M$ would imply $J_\lambda(t)>J_\lambda(u)$, a contradiction.
Put $s=f^{L-1}(t)$. Since $t\ge N(\lambda)\ge2$, one has $L\le s-1$,
so $f(s)<J_\lambda(t)<2f(s)<f(s+1)$ by Lemma~\ref{m:lem:growth}.
These intervals are disjoint and preserve order. If $W>Z=f(s)$ is another
endpoint, then $W\ge f(s+1)\ge2^{Z+1}$.
\end{proof}

\subsection{Nested common-head sums}
Fix any finite list of fixed nonzero limits $\lambda_1,\ldots,\lambda_l$.
Let $\mathcal A_0$ consist of natural constants, positive raw variables,
addition, and $f$, with arbitrary fixed nesting. Recursively set
\begin{equation}\label{m:eq:common-class}
 \mathcal A_r=\left\{c+\sum_i a_i h_{\lambda_r}(B_i):
 c\ge0,\ a_i>0,\ B_i\in\mathcal A_{r-1}\right\}.
\end{equation}
The sum can be empty. Each input supplies this finite level certificate.
There is no open remainder outside the head sum at a positive level.
For example $x+h_\lambda(y)$ does not pass that positive-level certificate
merely because both $x,y$ are allowed at level zero. Heads at different
levels may be completely different, and may exceed $\omega^\omega$.

\begin{theorem}[Common-head layers]\label{m:thm:common}
For $E,F\in\mathcal A_r$, each pointwise relation $E<F$, $E=F$, $E>F$
has an effectively constructed formula in $(\N;0,1,+,<,f)$ with the same
raw free variables and exactly the same truth at every positive valuation.
Hence arbitrary Boolean and quantified combinations of certified comparisons
are decidable.
\end{theorem}
\begin{proof}
Subtract the two expressions and write the row as
\[
 Q=K+\sum_{i<m}a_i J_{\lambda_r}(t_i),\qquad
 t_i=N(\lambda_r)+B_i.
\]
Coefficients and $K$ are signed integers. With no remaining atoms the sign
is the computable sign of $K$. Otherwise set
$C=\sum_i|a_i|$ and $L=\max(4,2C,4|K|+4)$.
On the branch with all $t_i\ge L$, partition the inputs by equal values,
combine coefficients, and delete zero blocks. If every block cancels,
$Q=K$. Otherwise its sign is the coefficient of the greatest surviving
input. Indeed let $W$ be its endpoint and $Z$ the greatest lower endpoint,
if any. Lemma~\ref{m:lem:endpoints} gives
$W\ge2^{Z+1}\ge2Z^2\ge4CZ$, since $Z\ge L\ge\max(4,2C)$.
All lower contributions have total absolute value less than
$2CZ\le W/2$. Also $|K|<W/4$, whereas the leading contribution has
absolute value greater than $W$. Its sign cannot change. Larger canceled
blocks play no part in this estimate. Enumerating the finitely many total
preorders of the labeled inputs produces a finite formula, guarded by
$t_i=t_j$, $t_i<t_j$, and $t_i\ge L$.

For the other branch choose a small remaining index, enumerate
$N(\lambda_r)\le c<L$, retain its exact guard $t_i=c$, evaluate the
closed $J_{\lambda_r}(c)$, fold $a_iJ_{\lambda_r}(c)$ into $K$, and
remove that occurrence. Recompute $L$ and repeat. Each small call removes
one occurrence; its possibly enormous updated constant is still a fixed
computable integer. Overlapping branches are harmless in a disjunction.
No open value is folded into a constant before its whole argument is fixed.

Induct now on $r$. At level zero comparisons are base-language atoms.
At higher levels every guard compares preceding-level expressions with
each other or with constants; a proposed constant below $N(\lambda_r)$
gives a constant truth value. Thus every generated predicate is covered
by the induction hypothesis. The lexicographic pair (level, remaining
head occurrences) proves termination: small folding decreases the second,
guard translation the first. Shared variables remain shared throughout.
The result is an exact base formula, to which
Theorem~\ref{m:thm:snyatkov} applies after any desired raw quantifiers are
added.
If cancellation removed a free name syntactically, conjoin the tautology
$x=x$ for that original name.  This gives exactly the original free-name
set without changing the predicate.
\end{proof}

For quantified combinations, translate the comparison atoms in place and
preserve the original Boolean structure.  A positive existential binder
becomes $\exists x\,(x\ge1\wedge\varphi)$, and a positive universal
binder becomes $\forall x\,(x\ge1\mathbin{\to}\varphi)$.  Bound names
can first be renamed to avoid capture.  Each atom may have its own
genuine finite certificate; its translation creates no independent raw
name for a derived value.  This is a parameter-preserving reduction to
the full base theory, so alternating binders are allowed here.

\begin{example}[Cancellation changes the controlling input]
Consider $2h_\lambda(x)+h_\lambda(y)-3h_\lambda(z)$. On the computed
large-input branch, a unique largest argument determines the sign by its
coefficient. But on the guard $x=z$, the row becomes exactly
$h_\lambda(y)-h_\lambda(x)$, whose sign is that of $y-x$; when all three
arguments agree the row is zero. The algorithm retains those guards rather
than treating three head values as independent large numbers. Small inputs
are evaluated by the finite folding procedure, so this is an exact
comparison theorem rather than only an eventual rule.
\end{example}

\subsection{Different heads in one control interval}
\begin{lemma}[Head order in a strict interval]\label{m:lem:head-window}
Let $\lambda>\mu$ be nonzero fixed limits. If
$t\ge\max(N(\lambda),N(\mu)+2)$ and
$N(\mu)\le u<f(t)$, then $J_\lambda(t)>J_\mu(u)$ and their endpoints
differ. Consequently, if finitely many norm inputs $t_i$ have minimum
$s\ge\max_iN(\lambda_i)+2$ and $\max_i t_i<f(s)$, their outputs are
ordered lexicographically by (ordinal head, norm input). Equality requires
equality of both, and different output blocks have distinct sparse endpoints.
\end{lemma}
\begin{proof}
The least limit above $\mu$ is $\mu+\omega$, of norm $N(\mu)+2$.
Controlled strictness at common norm $t$, or equality of the two heads,
gives
\[
 J_\lambda(t)\ge J_{\mu+\omega}(t)
       =J_\mu(f(t))+1>J_\mu(u).
\]
The exact middle identity follows by removing the final $\omega$ block:
its first predecessor retains $\mu$, assigns the remaining full budget
to its finite tail, and has total norm $f(t)$. If the two displayed
outputs had equal endpoints, their complete counts would satisfy
$L>M$; canceling the common $f$ iterates would give
$u=f^{L-M}(t)\ge f(t)$, a contradiction. For the finite collection, the
lower input is below $f(s)\le f(t)$ for every higher-head input $t$.
Apply this argument to different heads and
Lemma~\ref{m:lem:endpoints} within the same head.
\end{proof}

The interval inequality is strict for a reason. At $u=f(t)$,
\[
 J_{2\omega}(t)=f^2(t)+2,\qquad
 J_\omega(u)=f^2(t)+1.
\]
Their endpoints coincide, and
$J_{2\omega}(t)-2J_\omega(u)=-f^2(t)<0$ despite the positive coefficient
of the greatest ordinal head. An algorithm using head order at that
boundary would therefore be incorrect.

\subsection{Balanced heterogeneous layers}
The next certificate permits different sibling heads and an open lower-level
remainder. At level zero use the same $\mathcal A_0$ as before.  Define a
certificate recursively.  At a positive level it consists of one
preceding certificate $C$ and a finite ordered basis
$B_1,\ldots,B_k$, all certified by that same $C$.  The basis may have
repeated or dependent entries.  Expressions with the new certificate
are nonnegative affine combinations of the $B_j$ and atoms
\begin{equation}\label{m:eq:balanced-atoms}
 h_\lambda\!\left(c+\sum_j a_jB_j\right),
       \qquad c\ge0,\quad a_j\ge1.
\end{equation}
Each atom may have its own arbitrary fixed head and coefficients. The
essential condition is that \emph{every head argument includes every
common basis entry positively}. The basis is checked on the whole pair
being compared, including its outside remainder. A variable appearing
only in that remainder is not admitted by ignoring it while choosing the
basis. Zero multiples are discarded when computing actual supports.
For each fixed certificate its expression family is closed under
nonnegative affine combinations: collect constants, outside coefficients,
and atom lists.  Induction gives the same closure at the preceding
certificate.  Consequently the lower affine guards used below really do
belong to one preceding comparison problem; separate certificates are
not silently combined into an independent basis.

\begin{theorem}[Balanced comparison]\label{m:thm:balanced}
For two expressions with one finite balanced certificate, their three-way
pointwise comparison has an effectively constructed base formula in
$(\N;0,1,+,<,f)$ with exactly the same free variables. Arbitrary Boolean
and quantified combinations of such comparisons are decidable.
\end{theorem}
\begin{proof}
A signed comparison row has the form
\begin{equation}\label{m:eq:balanced-row}
 Q=K+\sum_jr_jB_j+\sum_{i<m}q_iJ_{\lambda_i}(t_i),\qquad
 t_i=b_i+\sum_j a_{ij}B_j,\quad b_i=N(\lambda_i)+c_i.
\end{equation}
The head-free case is already a signed affine comparison at the preceding
level. Otherwise put
\[
 S=\sum_jB_j,\quad C=\sum_i|q_i|,\quad C_0=\max_j|r_j|,
 \quad A=\max_{i,j}a_{ij},\quad B=\max_i b_i,
\]
with $C_0=0$ for an absent remainder and $A=1$ for an empty basis. Set
\begin{equation}\label{m:eq:balanced-L}
 L=\max\{16,2C,4C_0+4|K|+4,A+B+2,\max_iN(\lambda_i)+2\}.
\end{equation}
On the branch $s=\min_i t_i\ge L$ one has
\[
 S\le t_i\le AS+B,\qquad
 \max_i t_i\le As+B<s^2\le2^s\le f(s).
\]
Thus Lemma~\ref{m:lem:head-window} orders all head outputs by head and then
input. Partition the inputs within each equal head, combine coefficients,
and discard zero blocks. If all cancel, the remaining row is exactly the
preceding-level $K+\sum r_jB_j$.

If a block survives, let $W$ be the endpoint of the last one and $v$ its
norm input. Its coefficient has absolute value at least one, so its
contribution has magnitude greater than $W$. The sparse-endpoint estimate
bounds all lower head contributions by $W/2$, exactly as in the proof of
Theorem~\ref{m:thm:common}. The outside remainder is bounded using this
\emph{surviving} input, not a possibly canceled maximum:
\[
 \left|K+\sum_jr_jB_j\right|\le |K|+C_0v
          <v^2/4\le f(v)/4\le W/4.
\]
The surviving last coefficient therefore determines the strict sign.
The finitely many guarded preorders within equal heads express this rule
by comparisons of affine preceding-level expressions.

We must retain every small exception. Every positive expression in the
certificate satisfies $B(\vec x)\ge x$ for each raw variable in its
actual support. Induct through addition and $f$, and at a head use
$h_\lambda(e)>e$: the finite ordinal $e$ is below $\lambda+e$ and has
norm $e\le f(N(\lambda)+e)$, so controlled strictness and $\psi(e)=e$
apply. Fixed positive multiples preserve the property. If some $t_i<L$,
positivity of every $a_{ij}$ implies $S<L$. Therefore every raw variable
anywhere in the basis, including the outside remainder, is less than $L$.
Enumerate the finite box $[1,L-1]^V$ on its \emph{actual} support $V$;
evaluate each closed row and its inputs exactly, retaining precisely
points where some $t_i<L$ and the requested sign holds. Variables outside
$V$ remain unrestricted; empty support gives a single empty valuation.

All large-branch guards, and the all-canceled remainder, are affine
comparisons of the $B_j$ and thus have their preceding-level certificate.
The finite-point exception formula is already in the base language.
Induction on certificate level proves the claimed translation and
termination. Theorem~\ref{m:thm:snyatkov} supplies its arithmetic decision.
Repeated basis entries and shared variables never become independent
parameters in this construction.
Conjoining $x=x$ for any canceled original name makes the free-name
set exact.  The in-place Boolean and positive-binder translation following
Theorem~\ref{m:thm:common} then gives arbitrary quantified decisions.
\end{proof}

\begin{example}[A lower-head argument may be larger]
Let $E,F$ have a common preceding-level certificate and $\lambda>\mu$.
The row
\[
 3h_\lambda(E+F+2)-2h_\mu(2E+F+1)-F
\]
is balanced on basis $(E,F)$. Its two norm arguments need not have the
same order as their ordinal heads, and $F$ remains open outside them.
Nevertheless, once both norms exceed \eqref{m:eq:balanced-L}, the row is
positive by the higher-head coefficient. Its exact comparison is completed
by the finite support box. Omitting $F$ from one head argument would
remove the inference that a small input bounds the entire open remainder.
\end{example}

\subsection{One arbitrary core: direct Presburger comparison}
An especially useful specialization no longer requires a certificate for the
interior core. Let $E$ be any expression made from natural constants,
positive raw variables, addition, fixed nonnegative multiples, and arbitrary
fixed nonzero-limit functions $h_\lambda$, with arbitrary finite nesting.
No subtraction or exposed $f$ constructor occurs inside this stated grammar.
The only property needed below is $E\ge x$ for every variable in its actual
support, proved by the positive-source induction just given. It need not
have a balanced-level certificate, nor range independently over all naturals.

\begin{theorem}[Arbitrary single-core rows]\label{m:thm:single-core}
For the signed row
\begin{equation}\label{m:eq:single-core}
 Q(\vec x)=K+rE(\vec x)+
       \sum_i q_i h_{\lambda_i}(a_iE(\vec x)+c_i),
 \quad a_i\ge1,\ c_i\ge0,
\end{equation}
with fixed integer $K,r,q_i$ and fixed nonzero limit heads, each sign
predicate $Q<0$, $Q=0$, $Q>0$ has an effectively computable finite
Presburger formula in the original raw variables. Thus universal comparison
of these rows needs no arithmetic decision theorem about $f$.
\end{theorem}
\begin{proof}
Collect identical triples $(\lambda_i,a_i,c_i)$ and discard their zero
coefficients. If any heads remain, use \eqref{m:eq:balanced-L} with the
single formal basis $E$, $C_0=|r|$, and $b_i=N(\lambda_i)+c_i$. Increase
that bound to
\[
 T=\max\bigl(L,1+\max_{i,j}|b_i-b_j|\bigr).
\]
For $E\ge T$, the norms lie in one strict $f$ interval. Within equal
heads their affine inputs are ordered lexicographically by $(a_i,b_i)$:
a difference of integer slopes dominates the fixed offset difference.
Equal triples have already been collected. Lemma~\ref{m:lem:head-window}
and the remainder estimate in the balanced proof therefore show that
the sign is the coefficient of the greatest triple
$(\lambda_i,a_i,b_i)$. This use of the balanced \emph{numerical estimate}
requires no certificate or comparison algorithm for $E$ itself.
If no heads remain, use $T=\max(1,|K|+1)$; the final sign of $rE+K$ is
$\sgn(r)$ when $r\ne0$, otherwise $\sgn(K)$. Denote the computed final
sign by $\sigma$.

Let $V$ be the actual raw support of $E$. Empty support is handled by one
closed evaluation. Otherwise, outside $[1,T-1]^V$ there is an $x\ge T$,
and positivity gives $E\ge T$. Evaluate every row in the finite box exactly
and record its sign. For a desired sign $s$ the exact formula is
\[
 \left(\bigvee_{x\in V}x\ge T\right)\wedge[\sigma=s]
 \quad\vee\quad
 \bigvee_{\substack{v\in[1,T-1]^V\\\sgn Q(v)=s}}
       \bigwedge_{x\in V}x=v(x).
\]
It is a finite Presburger formula on the original variables. The box is
effective and each closed evaluation terminates. If the tail falsifies a
universal comparison, a concrete admissible countervaluation sets one
supported variable to $T$ and all others to one; the dominance property
certifies it without evaluating a giant collapse value.
\end{proof}

For example, even a core $E=h_\alpha(x)+h_\beta(h_\gamma(y)+z)$ outside
a proposed balanced certificate may be used in
$3h_\lambda(E+c)-2h_\mu(2E+d)-E$. Its support is still $(x,y,z)$.
This result is a comparison theorem for that row, not a graph definition
or a general entry-map algorithm for the core. An unrelated open term
outside $E$, a variable ordinal head, or a nonlinear dependence on $E$
is outside the statement.

